\chapter{Time series and temporal data}
\label{ch:temporal-data}

\begin{quote}
\emph{``Time is the most important feature in a time series --- and the most often mishandled.''}
\end{quote}

% ============================================================
\section{Datetime parsing}

The first step with temporal data is converting strings to proper datetime objects. Pandas provides robust parsing, but real-world date formats are surprisingly inconsistent.

\begin{minted}{python}
import pandas as pd
import numpy as np

# Jena Climate dataset: 10-minute weather measurements
# Download from: https://www.kaggle.com/datasets/mnassrib/jena-climate
df = pd.read_csv("jena_climate_2009_2016.csv")
print(df.head())
print(f"\nDate column dtype: {df['Date Time'].dtype}")
\end{minted}

\begin{minted}{python}
# Parse datetime
df["datetime"] = pd.to_datetime(df["Date Time"], format="%d.%m.%Y %H:%M:%S")
df = df.set_index("datetime").drop(columns=["Date Time"])

print(f"Date range: {df.index.min()} to {df.index.max()}")
print(f"Frequency: ~{(df.index[1] - df.index[0])}")
print(f"Total records: {len(df):,}")
\end{minted}

\begin{warningbox}
Date parsing ambiguity: is ``01/02/2024'' January 2 (US) or February 1 (Europe)? Always specify the format explicitly with the \texttt{format} parameter. Never rely on automatic inference for production code.
\end{warningbox}

\subsection{Handling multiple date formats}

\begin{minted}{python}
# When a column has mixed date formats
messy_dates = pd.Series(["2024-01-15", "15/01/2024", "Jan 15, 2024",
                         "15-Jan-2024", "20240115"])

# pd.to_datetime with infer_datetime_format handles many cases
parsed = pd.to_datetime(messy_dates, format="mixed", dayfirst=True)
print(parsed)

# For truly messy data, use a custom parser
def parse_flexible(date_str):
    """Try multiple date formats."""
    formats = ["%Y-%m-%d", "%d/%m/%Y", "%b %d, %Y",
               "%d-%b-%Y", "%Y%m%d"]
    for fmt in formats:
        try:
            return pd.to_datetime(date_str, format=fmt)
        except (ValueError, TypeError):
            continue
    return pd.NaT
\end{minted}

% ============================================================
\section{Resampling}

\begin{minted}{python}
# Downsample: 10-minute data to hourly
hourly = df.resample("h").mean()
print(f"10-min records: {len(df):,}")
print(f"Hourly records: {len(hourly):,}")

# Downsample to daily
daily = df.resample("D").agg({
    "T (degC)": ["mean", "min", "max"],
    "p (mbar)": "mean",
    "rh (%)": "mean",
})
daily.columns = ["temp_mean", "temp_min", "temp_max",
                 "pressure_mean", "humidity_mean"]
print(daily.head())
\end{minted}

\begin{minted}{python}
# Upsample: daily to hourly (with interpolation)
daily_sample = daily.head(30)
hourly_upsampled = daily_sample.resample("h").interpolate(method="linear")
print(f"Daily records: {len(daily_sample)}")
print(f"Hourly (upsampled): {len(hourly_upsampled)}")
\end{minted}

\begin{preprocesstip}
When downsampling, think carefully about the aggregation function. Temperature should use mean, precipitation should use sum, and minimum/maximum readings should use min/max respectively. Using the wrong aggregation produces meaningless results.
\end{preprocesstip}

% ============================================================
\section{Lag features}

\begin{minted}{python}
# Lag features: use past values as predictors
daily["temp_lag1"] = daily["temp_mean"].shift(1)   # yesterday
daily["temp_lag7"] = daily["temp_mean"].shift(7)   # last week
daily["temp_lag30"] = daily["temp_mean"].shift(30)  # last month

# Change features
daily["temp_change_1d"] = daily["temp_mean"].diff(1)
daily["temp_change_7d"] = daily["temp_mean"].diff(7)

# Percentage change
daily["temp_pct_change"] = daily["temp_mean"].pct_change()

print(daily[["temp_mean", "temp_lag1", "temp_lag7",
             "temp_change_1d"]].head(10))
\end{minted}

\begin{minted}{python}
# Autocorrelation: check if lag features are useful
import matplotlib.pyplot as plt
from pandas.plotting import autocorrelation_plot

fig, axes = plt.subplots(1, 2, figsize=(14, 5))

# Autocorrelation plot
autocorrelation_plot(daily["temp_mean"].dropna(), ax=axes[0])
axes[0].set_title("Temperature autocorrelation")
axes[0].set_xlim(0, 400)

# Partial autocorrelation (useful for determining lag order)
from statsmodels.graphics.tsaplots import plot_pacf
plot_pacf(daily["temp_mean"].dropna(), lags=50, ax=axes[1])
axes[1].set_title("Partial autocorrelation")

plt.tight_layout()
plt.savefig("autocorrelation.png", dpi=150)
plt.show()
\end{minted}

\begin{warningbox}
Lag features introduce NaN values at the beginning of the series. A lag-30 feature has 30 missing values. Always account for this by either dropping the initial rows or using forward-fill imputation.
\end{warningbox}

% ============================================================
\section{Rolling statistics}

\begin{minted}{python}
# Rolling (moving) statistics smooth out noise
daily["temp_rolling7"] = daily["temp_mean"].rolling(window=7).mean()
daily["temp_rolling30"] = daily["temp_mean"].rolling(window=30).mean()
daily["temp_rolling_std7"] = daily["temp_mean"].rolling(window=7).std()

# Exponentially weighted moving average (more recent = more weight)
daily["temp_ewm7"] = daily["temp_mean"].ewm(span=7).mean()

print(daily[["temp_mean", "temp_rolling7", "temp_rolling30",
             "temp_ewm7"]].tail(10))
\end{minted}

\begin{minted}{python}
# Visualize rolling statistics
fig, ax = plt.subplots(figsize=(14, 6))

daily["temp_mean"].plot(ax=ax, alpha=0.3, label="Daily mean")
daily["temp_rolling7"].plot(ax=ax, label="7-day rolling mean")
daily["temp_rolling30"].plot(ax=ax, label="30-day rolling mean",
                             linewidth=2)

ax.set_ylabel("Temperature (degC)")
ax.set_title("Jena Climate: Temperature with rolling averages")
ax.legend()
plt.tight_layout()
plt.savefig("rolling_stats.png", dpi=150)
plt.show()
\end{minted}

% ============================================================
\section{Trend and seasonality decomposition}

\begin{minted}{python}
from statsmodels.tsa.seasonal import seasonal_decompose

# Decompose into trend + seasonal + residual
# Use monthly data for cleaner decomposition
monthly = daily["temp_mean"].resample("ME").mean()

decomposition = seasonal_decompose(monthly, model="additive", period=12)

fig = decomposition.plot()
fig.set_size_inches(14, 10)
plt.tight_layout()
plt.savefig("decomposition.png", dpi=150)
plt.show()
\end{minted}

\begin{minted}{python}
# Extract components as features
monthly_df = pd.DataFrame({
    "observed": decomposition.observed,
    "trend": decomposition.trend,
    "seasonal": decomposition.seasonal,
    "residual": decomposition.resid,
})
print(monthly_df.head(15))
\end{minted}

\begin{preprocesstip}
The decomposition residual is what remains after removing trend and seasonality. A large residual indicates an unusual observation --- it can serve as an outlier detection mechanism for time series.
\end{preprocesstip}

% ============================================================
\section{Air Quality case study}

\begin{minted}{python}
# Air Quality UCI dataset: hourly pollutant measurements
df_air = pd.read_csv("AirQualityUCI.csv", sep=";", decimal=",",
                      parse_dates={"datetime": ["Date", "Time"]},
                      dayfirst=True)

# Replace sentinel values
df_air = df_air.replace(-200, np.nan)

# Select key columns
cols = ["datetime", "CO(GT)", "NO2(GT)", "T", "RH"]
df_air = df_air[cols].set_index("datetime").dropna(how="all")

# Resample to daily
air_daily = df_air.resample("D").mean()

# Create temporal features
air_daily["hour_of_day_peak"] = df_air["CO(GT)"].resample("D").idxmax().dt.hour
air_daily["co_lag1"] = air_daily["CO(GT)"].shift(1)
air_daily["co_rolling7"] = air_daily["CO(GT)"].rolling(7).mean()
air_daily["co_weekend"] = air_daily.index.dayofweek.isin([5, 6]).astype(int)

print(air_daily.head(10))
\end{minted}

% ============================================================
\section{Datetime feature extraction summary}

\begin{minted}{python}
def extract_datetime_features(df, date_col):
    """Extract a comprehensive set of datetime features."""
    dt = df[date_col] if date_col in df.columns else df.index
    features = pd.DataFrame(index=df.index)

    features["year"] = dt.year
    features["month"] = dt.month
    features["day"] = dt.day
    features["dayofweek"] = dt.dayofweek
    features["hour"] = dt.hour if hasattr(dt, "hour") else 0
    features["quarter"] = dt.quarter
    features["is_weekend"] = dt.dayofweek.isin([5, 6]).astype(int)
    features["day_of_year"] = dt.dayofyear
    features["week_of_year"] = dt.isocalendar().week.astype(int)

    # Cyclical encoding
    features["month_sin"] = np.sin(2 * np.pi * dt.month / 12)
    features["month_cos"] = np.cos(2 * np.pi * dt.month / 12)
    features["dow_sin"] = np.sin(2 * np.pi * dt.dayofweek / 7)
    features["dow_cos"] = np.cos(2 * np.pi * dt.dayofweek / 7)

    return features
\end{minted}

% ============================================================
\section{Exercises}

\begin{exercise}
\begin{enumerate}
  \item Load the Jena Climate dataset. Parse the datetime column and resample to daily frequency. Plot the daily mean temperature for 2015.
  \item Create lag features (1, 7, 14, 30 days) for temperature. Compute the correlation of each lag feature with the current temperature. Which lag has the highest correlation?
  \item Apply seasonal decomposition to the Air Quality CO data. Is there a clear seasonal pattern? What does the trend tell you?
  \item Build a comprehensive feature extraction function for the Jena Climate dataset that creates: lag features, rolling statistics, cyclical time encodings, and trend indicators. Apply it and report the resulting number of features.
  \item Using the Air Quality dataset, investigate whether CO levels are significantly different on weekdays vs.\ weekends. Create appropriate visualizations.
\end{enumerate}
\end{exercise}

% ============================================================
\section{Chapter summary}

\begin{itemize}
  \item Always parse dates explicitly with a known format; never rely on automatic inference in production.
  \item Resampling changes temporal resolution: downsampling aggregates, upsampling interpolates.
  \item Lag features capture temporal dependencies: yesterday's temperature predicts today's.
  \item Rolling statistics (moving average, moving std) smooth noise and reveal trends.
  \item Seasonal decomposition separates trend, seasonality, and residuals --- each can serve as a feature.
  \item Cyclical encoding (sine/cosine) ensures that December is close to January in feature space.
\end{itemize}
